% year: תשפ"ו
% type: בוחן
% semester: א
% official_solution: yes
% source: exams/בחנים/תשפ״ו, בוחן, 9141.pdf
% part: א
% number: 1
% points: 20
% categories: functions

\begin{question}
נגדיר פונקציה על ידי
\[
f(x)=\begin{cases}
x^3 & x<1,\\
2x+1 & x\ge 1
\end{cases}
\]
\begin{enumerate}
  \item (5 נק') שרטטו את גרף הפונקציה.
  \item (15 נק') האם פונקציה זו הפיכה כפונקציה מ-$\R$ ל-$\R$?
  אם כן - חשבו את הפונקציה ההופכית שלה $f^{-1}:\R\to\R$.
  אם לא - נמקו.
\end{enumerate}
\end{question}

\begin{hint}
פונקציה הפיכה מ-$\R$ ל-$\R$ חייבת להיות גם חח"ע וגם על $\R$. מהי התמונה של כל אחד משני החלקים?
\end{hint}

\begin{solution}
\begin{enumerate}
  \item \textbf{הגרף:} עבור $x<1$ זהו גרף הפרבולה הקובית $y=x^3$: עולה ממש, עובר בנקודות $(-1,-1)$, $(0,0)$ (נקודת פיתול עם משיק אופקי), ומתקרב לנקודה $(1,1)$ — נקודה זו \emph{אינה} שייכת לגרף (עיגול ריק). עבור $x\ge1$ זו קרן ישרה עם שיפוע 2 שמתחילה בנקודה $(1,3)$ (עיגול מלא) ועוברת ב-$(2,5)$. בנקודה $x=1$ יש "קפיצה" מגובה 1 לגובה 3.
  \item \textbf{הפונקציה אינה הפיכה כפונקציה מ-$\R$ ל-$\R$}, כי היא אינה על $\R$.

  נמצא את התמונה. עבור $x<1$: $x^3$ עולה ממש, ולכן $x^3<1^3=1$; כלומר ערכי החלק הראשון קטנים מ-1 (ולמעשה מכסים את $(-\infty,1)$). עבור $x\ge1$: $2x+1\ge3$; כלומר ערכי החלק השני גדולים או שווים ל-3 (ומכסים את $[3,\infty)$). לכן
  \[
  \operatorname{Im}f=(-\infty,1)\cup[3,\infty).
  \]
  בפרט, לא קיים $x\in\R$ עם $f(x)=2$: אם $x<1$ אז $f(x)<1$, ואם $x\ge1$ אז $f(x)\ge3$. לכן $f$ אינה על $\R$, ולכן אין לה הופכית $f^{-1}:\R\to\R$.

  (הערה: $f$ כן חח"ע — היא עולה ממש, שכן כל חלק עולה ממש וכל ערך של החלק הראשון קטן מכל ערך של החלק השני. לכן $f$ הפיכה כפונקציה מ-$\R$ אל התמונה שלה $(-\infty,1)\cup[3,\infty)$, עם $f^{-1}(y)=\sqrt[3]{y}$ עבור $y<1$ ו-$f^{-1}(y)=\frac{y-1}{2}$ עבור $y\ge3$; אבל לא כפונקציה מ-$\R$ ל-$\R$.)
\end{enumerate}
\end{solution}
